\section{Answers to the Research Questions}
\textbf{RQ1: Does classification vary over recording time, and can stable predictions be wrong?} Yes, in the historical Study A spectral SVM setting. First-minute balanced accuracy was 73.48\% for AD/CN and 73.72\% for FTD/CN, yet 44.35\% and 33.66\% of correlated participant-repeat evaluations changed labels across five minutes. Stable-wrong evaluations accounted for 8.70\% and 13.17\%. In the independent Isfahan MCI negative control, many decisions remained stable even when first-minute classification was poor. These are different models and datasets; neither proves EEGraft itself has the same five-minute behavior.

\textbf{RQ2: How does EEGraft compare with classical models and external data?} The custom frozen-CBraMod DSRF decoder reached 77.06\% balanced accuracy on 52 source development participants, compared with 76.61\% for XGBoost. On the original locked 13-person source holdout, DSRF reached 70.24\% versus 77.38\% for XGBoost. On the 57-recording locked CAUEEG broad-Dementia/Normal task, the original 0.50 thresholds yielded 68.38\% and 58.57\%, respectively. Neither paired model comparison established a general superiority advantage, and the external task is not an AD-specific replication.

\textbf{RQ3: Does the decision rule affect reliability interpretation?} Yes, on exploratory CAUEEG development condition pairs. At 0.50, XGBoost changed the eyes-open versus eyes-closed mean-score decision for 7.18\% of eligible training recordings; at the training-selected threshold 0.70, it changed for 33.17\%, without changing its model weights. Locked-test unanimity also revealed 12 wholly stable-but-wrong XGBoost recordings. Train-fitted probability mapping changed Brier scores and balanced accuracy, but these secondary analyses were not independently preregistered confirmation.

\textbf{RQ4: Which validity limits matter?} Age differences between CAUEEG Normal and Dementia recordings were large, while adding EEGraft scores to age-only validation improved AUROC only slightly, with an interval crossing zero. Eye-state score differences were entangled with recording time. Diagnostic labels shifted from AD/CN to broad Dementia/Normal; patient uniqueness was unverified for the restricted recording data. These factors limit medical or causal interpretation.

\section{Scientific and Engineering Contributions}
This thesis delivers a working frozen-CBraMod EEGraft decoder with a trainable dual-scale residual fusion head, documented architecture ablations, participant-disjoint development comparisons, and a separately reported 13-person locked source test. It also provides a frozen CAUEEG transfer comparison without target weight updates, plus an evaluation of incorrect unanimity, decision-threshold dependence, calibration and age/condition influences. The historical NEMAR temporal study and independent Isfahan MCI analysis show why stable labels must be checked against ground truth rather than treated as reliability in isolation.

The transferable result for future researchers is an evaluation procedure: define the unit of analysis, compare against classical baselines, preserve test isolation and original decision rules, report stable incorrect predictions, and check demographic or recording-context effects before interpreting apparent generalization. These are practical contributions supported by the specific experiments, not claims to invent the components or to establish a clinical biomarker.

\section{Recommended Future Work}
The next study should use an explicitly matched person-level protocol and compare the same frozen DSRF, XGBoost and SVM models across recording minutes. A larger external sample with verified unique patients, harmonized acquisition and independent clinical labels is needed before drawing reliable transfer conclusions. A prospective study should control or match for age and recording condition, and it should test thresholds and calibration rules chosen without using final evaluation outcomes. Multiple training seeds, cost and energy measurements, physically realistic artifact perturbations and source-audited EEG sample previews would strengthen the engineering claims. Controlled CAUEEG access and any publication permission should be documented before releasing restricted signals.

\section{Final Conclusion}
EEGraft demonstrates a testable dual-scale decoder and, more importantly, an evidence-based way to examine EEG predictions beyond one accuracy number. Across source, external and negative-control analyses, the work identifies stable incorrect decisions, threshold-driven changes in consistency and important cohort differences. Its strongest outcome is a reproducible reporting approach that helps future researchers distinguish a model that predicts consistently from one that predicts correctly, while keeping the evidence grounded in the cohorts actually evaluated.
